\documentclass[a4paper,11pt]{article}
%% Compiles both from this folder (pdflatex tp4.tex) and from the root of the project
%% (Overleaf: Menu > Main document = tp4_pacbayes_intro/tp4.tex).
\makeatletter
\IfFileExists{tp4_pacbayes_intro/tp4.tex}{\gdef\tpdir{tp4_pacbayes_intro/}\gdef\tpcommon{common/}}{\gdef\tpdir{}\gdef\tpcommon{../common/}}
\edef\input@path{{\tpdir}{\tpcommon}}
\makeatother
\input{tpsetup}
\renewcommand{\polytitle}{Statistical Supervised Machine Learning -- TP 4}
\graphicspath{{figures/}{tp4_pacbayes_intro/figures/}{../common/logos/}{common/logos/}}

\begin{document}
\sloppy
\sheethead{TP 4 -- From an error on a sample to a PAC-Bayesian guarantee}{%
Fourth of five practical sessions (1: protocol, 2: bagging, 3: boosting, 4 and 5: PAC-Bayes). This session follows
Sections~13 and~14 of the lecture notes, step by step. On a toy problem in which every true risk is known exactly, so that
we can draw thousands of samples and \emph{count} how often a bound fails, we (1) certify one classifier fixed in advance
with Hoeffding's inequality and with the sharper kl bound, (2) see what goes wrong when the classifier is chosen with the
sample and repair it with the union bound, (3) spread the confidence over the classifiers with a \emph{prior} (Occam's
bound), and then move from one classifier to a \emph{distribution} over classifiers: (4) the KL divergence, (5) the first
PAC-Bayesian bound of the notes, with its parameter $\lambda$, and the bounds of McAllester and Seeger used in practice,
(6) the Gibbs posterior, and (7) a check that all these bounds hold. Majority votes are the subject of TP~5. The session
lasts \textbf{three hours}; Part~\ref{p4:far} contains the difficult questions.}

\begin{sessionbox}[About this session -- TP 4: From an error on a sample to a PAC-Bayesian guarantee]
\textbf{Goal.} By the end of the session you should be able to compute the test-set and kl bounds of a fixed classifier,
explain why a selected classifier needs a union bound and how a prior spreads the confidence, compute the Gibbs risk and
the KL divergence of a posterior, compute the PAC-Bayesian bound with parameter $\lambda$ and the bounds of McAllester and
Seeger, explain the trade-off made by the Gibbs posterior, and check by simulation that the bounds hold.\par
\textbf{To reread before the session.} Section~13 (from an error on a sample to a guarantee: Corollary~13.3,
Theorem~13.5, Algorithm~12, Theorems~13.7 and~13.8) and Section~14 (the PAC-Bayesian bound: Definition~14.1,
Theorems~14.4 and~14.5, Proposition~14.6) of the lecture notes. The exercises of Sections~13 and~14 of the exercise
booklet are a good preparation.\par
\textbf{Timing (3 hours).}\quad
\begin{tabular}[t]{ll}Setting up and Part~1 & 35 min\\
Part~2 & 15 min\\
Part~3 & 20 min\\
Part~4 & 20 min\\
Part~5 & 35 min\\
Part~6 & 30 min\\
Part~7 & 25 min\\
Part~8 \hard & at home\end{tabular}\par
\textbf{To hand in.} Your notebook with the tables of Parts~1 to~7, the figures of Parts~1, 5 and~6, and your proofs.\par
\textbf{Tools.} \code{numpy} and \code{matplotlib} (and the helpers of \code{common/tptools.py}).
\end{sessionbox}

\medskip
\noindent\textbf{The toy problem.} The input $X$ is uniform on $[0,1]$. The clean label is $\sign(X-\theta^\star)$ with
$\theta^\star=0.37$, and it is flipped with probability $\eta=0.15$. The class $\Hcal$ contains the $K=102$ threshold
classifiers $h_{t,s}(x)=s\cdot\sign(x-t)$, $t\in\{0,0.02,\dots,1\}$, $s\in\{-1,+1\}$. Unless stated otherwise,
$\delta=0.05$.

\begin{enumerate}
\item Show that $\RD(h_{t,+1})=\eta+(1-2\eta)|t-\theta^\star|$ and $\RD(h_{t,-1})=1-\RD(h_{t,+1})$.\\
      What is the best classifier of $\Hcal$?
\item Write \code{sample(m, rng)} and a \emph{vectorized} function \code{emp\_risks(x, y)} returning the $102$
      empirical risks (build the $102\times m$ matrix of predictions).
\end{enumerate}

%% =====================================================================
\tppart{One classifier fixed in advance: Hoeffding and kl bounds (Sections 13.1--13.4)}\label{p4:fixed}

Fix $h=h_{0.6,+1}$, whose risk is $p=\RD(h)=0.311$. For a sample of size $m$, $m\RS(h)\sim\mathcal B(m,p)$: you can
simulate $\RS(h)$ directly with \code{rng.binomial}.

\begin{enumerate}
\item \textbf{Test-set bound.} From Hoeffding's inequality (Theorem~13.2),
      $\P\big(\RD(h)-\RS(h)\geq\varepsilon\big)\leq e^{-2m\varepsilon^2}$, prove Corollary~13.3: with probability at least
      $1-\delta$, $\RD(h)\leq\RS(h)+\sqrt{\ln(1/\delta)/(2m)}$.\\
      Why must $h$ be fixed before the sample is drawn?
\item \textbf{kl bound.} Theorem~13.5 gives, with probability at least $1-\delta$,
      $\RD(h)\leq\klup\big(\RS(h),\ln(1/\delta)/m\big)$, where
      $\klup(q,\varepsilon)=\max\{p\in[q,1]:\kl(q\|p)\leq\varepsilon\}$ and $\kl(q\|p)=q\ln\frac qp+(1-q)\ln\frac{1-q}{1-p}$
      (Definition~13.4).\\
      Implement $\klup$ by bisection (Algorithm~12), for whole arrays of $q$ at once, and check it against
      \code{kl\_inv\_upper} of \code{common/tptools.py}; check that $\klup(0.1,0.05)\approx0.220$ (example of
      Section~13.4).\\
      Using Pinsker's inequality $\kl(q\|p)\geq2(q-p)^2$, show that the kl bound is never worse than the test-set bound.
\item For $m\in\{50,100,500,2000\}$, draw $10\,000$ samples and compute the mean of each bound and its \emph{coverage}
      (the proportion of samples for which $\RD(h)$ is below the bound).\\
      Comment.
\item For $m=100$, plot the width ``bound $-\RS(h)$'' of both bounds as a function of the observed error
      $\RS(h)\in[0,0.5]$.\\
      When is the kl bound much better?
\end{enumerate}

%% =====================================================================
\tppart{What goes wrong when we choose: the union bound (Section 13.5)}\label{p4:union}

\begin{enumerate}
\item Using the union bound (Lemma~13.6), prove Theorem~13.7: with probability at least $1-\delta$, \emph{for all}
      $k\leq K$ simultaneously, $\RD(h_k)\leq\RS(h_k)+\sqrt{\ln(K/\delta)/(2m)}$; give its kl version
      $\RD(h_k)\leq\klup\big(\RS(h_k),\ln(K/\delta)/m\big)$.\\
      Why does it hold for a classifier \emph{selected} on the sample, whatever the selection rule?
\item For $m\in\{50,100,400\}$ and $5000$ samples, select the empirical risk minimizer (ERM)
      $\hat h$ among the $102$ thresholds.\\
      Compute the mean of $\RD(\hat h)-\RS(\hat h)$, and the coverage of (a) the
      test-set bound with $\delta$ (as if $\hat h$ were fixed), (b) the union bound for $\hat h$, (c) the union bound for
      \emph{all} the classifiers at once.\\
      Comment.
\end{enumerate}

%% =====================================================================
\tppart{Occam's bound: spreading the confidence with a prior (Section 13.6)}\label{p4:occam}

Let $P$ be a probability distribution on $\Hcal$ (the \emph{prior}), chosen \emph{before} seeing the data. The kl version
of Occam's bound (Theorem~13.8) states that, with probability at least $1-\delta$,
\[
\forall h\in\Hcal,\qquad \RD(h)\leq\klup\Big(\RS(h),\ \frac{\ln\frac1{P(h)}+\ln\frac1\delta}{m}\Big).
\]

\begin{enumerate}
\item Prove it, by applying the kl bound of Part~\ref{p4:fixed} to each $h$ with the confidence $\delta P(h)$ and the union
      bound.\\
      What do you recover with the uniform prior?
\item On the toy problem ($m=100$, $2000$ samples), compare three priors: uniform; ``informative'', with weights
      proportional to $e^{-(t-0.4)^2/(2\times0.1^2)}$ on the threshold and $90\%$ of the mass on $s=+1$; ``misleading'',
      the same centred on $0.9$.\\
      For each prior, report the mean bound and the mean true risk of the ERM and of the classifier that minimizes the
      \emph{bound}, and the coverage for all $h$ at once.
\item Can you choose the prior after looking at the data?\\
      What happens if the prior is bad?
\end{enumerate}

%% =====================================================================
\medskip
\noindent\textbf{From one classifier to a distribution over classifiers.} We keep the toy problem and its class $\Hcal$ of
$K=102$ thresholds; the prior $P$ is now uniform on $\Hcal$. A \emph{posterior} is any distribution $Q$ on $\Hcal$,
possibly chosen using the sample $S$ of size $m$. We recall (Sections~12 and~14 of the lecture notes):
\begin{itemize}
\item the \textbf{Gibbs classifier} $\GQ$ draws $h\sim Q$ and predicts $h(x)$; its risks are $\RD(\GQ)=\sum_hQ(h)\RD(h)$ and
      $\RS(\GQ)=\sum_hQ(h)\RS(h)$ (Definition~12.1); the \textbf{majority vote} $\BQ$ is the subject of TP~5;
\item the \textbf{KL divergence} $\KL(Q\|P)=\sum_hQ(h)\ln\frac{Q(h)}{P(h)}$, with $0\ln0=0$ (Definition~14.1);
\item the \textbf{PAC-Bayesian bound} of Theorem~14.4: for $\lambda>0$ \emph{fixed before seeing $S$}, with probability at
      least $1-\delta$, for all $Q$ simultaneously,
      \[
      \RD(\GQ)\leq\RS(\GQ)+\frac{\KL(Q\|P)+\ln\frac1\delta}{\lambda}+\frac{\lambda}{8m};
      \]
\item the bounds of \textbf{McAllester} and \textbf{Seeger} (Theorem~14.5): for $m\geq8$, with probability at least
      $1-\delta$, for all $Q$ simultaneously,
      \begin{gather*}
      \text{(McAllester)}\quad\RD(\GQ)\leq\RS(\GQ)+\sqrt{\frac{\KL(Q\|P)+\ln\frac{2\sqrt m}{\delta}}{2m}},\\
      \text{(Seeger)}\quad\RD(\GQ)\leq\klup\Big(\RS(\GQ),\ \frac{\KL(Q\|P)+\ln\frac{2\sqrt m}{\delta}}{m}\Big).
      \end{gather*}
\end{itemize}

\begin{tip}
Work with \emph{vectors}: a posterior is an array of $102$ nonnegative numbers summing to $1$, so that
$\RS(\GQ)$ is \code{Q @ r} and $\RD(\GQ)$ is \code{Q @ R}. To compute $\KL(Q\|P)$, sum only over the $h$ with $Q(h)>0$.
To compute weights $\propto e^{a_h}$ for large $a_h$, subtract $\max_ha_h$ before taking the exponential.
\end{tip}

%% =====================================================================
\tppart{Posteriors, Gibbs risk and KL divergence (Sections 12.2 and 14.1)}\label{p4:post}

\begin{enumerate}
\item Show that, for the uniform prior, $\KL(Q\|P)=\ln K-H(Q)$ with $H(Q)=-\sum_hQ(h)\ln Q(h)\geq0$, so that
      $0\leq\KL(Q\|P)\leq\ln K$.\\
      What is the KL of the uniform distribution on $k$ of the $K$ classifiers (Example~14.1)? Of a Dirac mass?
\item Draw one sample of size $m=100$.\\
      For the following posteriors, compute $\RS(\GQ)$, $\KL(Q\|P)$ and the true Gibbs risk $\RD(\GQ)$:
      (a) the Dirac mass on the ERM $\hat h=h_{\hat t,\hat s}$; (b) the prior itself; (c) the uniform distribution on the $5$
      classifiers with the smallest empirical risks; (d) the uniform distribution on the $h_{t,+1}$ with $|t-\hat t|\leq0.1$.\\
      Which posterior would you like to certify? What makes it ``complex''?
\end{enumerate}

%% =====================================================================
\tppart{The PAC-Bayesian bounds (Sections 14.3--14.5)}\label{p4:bounds}

\begin{enumerate}
\item On the sample of Part~\ref{p4:post}, compute the bound of Theorem~14.4 for the four posteriors and
      $\lambda\in\{10,30,100,300\}$.\\
      Plot it as a function of $\lambda\in[1,10^4]$ for the posteriors (a) and (d).\\
      Which term dominates for a small $\lambda$, for a large $\lambda$?
\item For a fixed $Q$, show that the value of $\lambda$ that minimizes the bound is
      $\lambda^\star=\sqrt{8m\big(\KL(Q\|P)+\ln\frac1\delta\big)}$, and that the minimum is
      $\RS(\GQ)+\sqrt{\big(\KL(Q\|P)+\ln\frac1\delta\big)/(2m)}$.\\
      Why can we \emph{not} use $\lambda^\star$, computed from the data, in Theorem~14.4?\\
      Show that choosing $\lambda$ in a grid of $k$ values fixed in advance is allowed if $\ln\frac1\delta$ is replaced by
      $\ln\frac k\delta$, and compute this certificate for the four posteriors with the grid $\{10,30,100,300\}$.
\item Compute McAllester's and Seeger's bounds for the four posteriors, and compare all the certificates.\\
      Compare the certificate of the Dirac mass with Occam's bound for $\hat h$ (Part~\ref{p4:occam}).\\
      Why is the PAC-Bayesian bound a little looser here?
\end{enumerate}

%% =====================================================================
\tppart{The Gibbs posterior (Section 14.6)}\label{p4:gibbs}

For $\lambda>0$, the \emph{Gibbs posterior} is $Q_\lambda(h)=\frac{P(h)\,e^{-\lambda\RS(h)}}{Z_\lambda}$
(Proposition~14.6): it minimizes $\RS(\GQ)+\frac{\KL(Q\|P)}{\lambda}$, hence the bound of Theorem~14.4 for this $\lambda$.

\begin{enumerate}
\item What is $Q_\lambda$ for $\lambda=0$ (where the formula still makes sense), and what is the limit of $Q_\lambda$ when $\lambda\to\infty$?
\item On the sample of Part~\ref{p4:post}, for $\lambda\in\{0\}\cup[1,10^4]$ (logarithmic grid of $41$ values), plot
      $\RS(G_{Q_\lambda})$, $\KL(Q_\lambda\|P)/m$, Seeger's bound and $\RD(G_{Q_\lambda})$ against $\lambda$.\\
      Plot $Q_\lambda$ for a few values of $\lambda$.
\item Select $\lambda$ by minimizing Seeger's bound.\\
      Why is Seeger's bound still valid for the selected posterior, without any union bound over the grid (compare with
      Part~\ref{p4:union} and with question~2 of Part~\ref{p4:bounds})?
\end{enumerate}

%% =====================================================================
\tppart{Do the bounds hold?}\label{p4:check}

\begin{enumerate}
\item For $m\in\{50,100,400\}$, draw $2000$ samples.\\
      On each, select $\lambda$ by minimizing Seeger's bound over the grid of Part~\ref{p4:gibbs}, and record
      $\RS(G_{Q_\lambda})$, $\RD(G_{Q_\lambda})$, Seeger's bound and whether it holds; record also the bound of
      Theorem~14.4 for the same posterior, with the union bound over the $41$ positive values of the grid.\\
      Compare with Occam's bound for the ERM and its true risk.
\item Comment on the coverage, on the tightness of the bounds, and on the comparison with the ERM.
\end{enumerate}

%% =====================================================================
\tppart{Going further \hard}\label{p4:far}

\begin{enumerate}
\item \textbf{The KL divergence is nonnegative.} Using Jensen's inequality for the concave function $\ln$, prove that
      $\KL(Q\|P)\geq0$, with equality if and only if $Q=P$ (Proposition~14.2).
\item \textbf{The Gibbs posterior is optimal.} Show that, for $\lambda>0$,
      $\lambda\RS(\GQ)+\KL(Q\|P)=\KL(Q\|Q_\lambda)-\ln Z_\lambda$, where
      $Z_\lambda=\sum_hP(h)e^{-\lambda\RS(h)}$.\\
      Deduce Proposition~14.6.
\item \textbf{Proof of Theorem~14.4.} Let $\Delta(h)=\RD(h)-\RS(h)$.\\
      (a) Using Hoeffding's lemma, $\E_S\big[e^{\lambda\Delta(h)}\big]\leq e^{\lambda^2/(8m)}$ for each fixed $h$, show that
      $\E_S\E_{h\sim P}\big[e^{\lambda\Delta(h)}\big]\leq e^{\lambda^2/(8m)}$.\\
      (b) Deduce with Markov's inequality (Lemma~13.1) that, with probability at least $1-\delta$,
      $\E_{h\sim P}\big[e^{\lambda\Delta(h)}\big]\leq\frac1\delta e^{\lambda^2/(8m)}$.\\
      (c) Conclude with the change of measure (Lemma~14.3),
      $\E_{h\sim Q}[\phi(h)]\leq\KL(Q\|P)+\ln\E_{h\sim P}\big[e^{\phi(h)}\big]$, applied to $\phi=\lambda\Delta$.\\
      Why does the result hold for all $Q$ simultaneously, and why must $\lambda$ be fixed in advance?
\end{enumerate}

\end{document}
